% Part: model-theory
% Chapter: basics
% Section: dlo

\documentclass[../../../include/open-logic-section]{subfiles}

\begin{document}

\olfileid{mod}{bas}{dlo}
\section{ترتیب‌های خطی چگال}

\begin{defn}
  یک \emph{ترتیب خطی چگال بدون نقاط انتهایی} !!a{structure} چون
  $\Struct{M}$ برای !!{language}ی شامل تنها یک
  !!{predicate}~$<$ دوموضعی است که جمله‌های زیر را برآورده می‌کند:
  \begin{enumerate}
  \item $\lforall[x][\lnot x < x]$؛
  \item $\lforall[x][\lforall[y][\lforall[z][(x < y \lif (y < z \lif x
    <z ))]]]$؛
  \item $\lforall[x][\lforall[y][(x< y \lor \eq[x][y] \lor y < x)]]$؛
  \item $\lforall[x][\lexists[y][x < y]]$؛
  \item $\lforall[x][\lexists[y][y < x]]$؛
  \item $\lforall[x][\lforall[y][(x < y \lif \lexists[z][(x < z \land
        z < y))]]]$.
 \end{enumerate}
\end{defn}

\begin{thm}\ollabel{thm:cantorQ}
  هر دو ترتیب خطی چگال بدون نقاط انتهایی که !!{enumerable} باشند
  همریخت‌اند.
\end{thm}

\begin{proof}
  بگذارید $\Struct{M_1}$ و~$\Struct{M_2}$ ترتیب‌های خطی چگال بدون نقاط
  انتهایی و !!{enumerable} باشند، به‌گونه‌ای که
  ${<_1} = \Assign{<}{M_1}$ و ${<_2} = \Assign{<}{M_2}$؛ و بگذارید
  $\PIso{I}$ مجموعهٔ همهٔ همریختی‌های جزئی میان آن‌ها باشد. $\PIso{I}$
  ناتهی است، زیرا دست‌کم $\emptyset \in \PIso{I}$. نشان می‌دهیم
  $\PIso{I}$ خاصیت پس‌وپیش را برآورده می‌کند. آنگاه
  $\Struct{M_1} \iso[p] \Struct{M_2}$، و قضیه بنا بر
  \olref[pis]{thm:p-isom1} نتیجه می‌شود.

  برای نشان‌دادن اینکه $\PIso{I}$ خاصیت پیش را برآورده می‌کند،
  $p \in \PIso{I}$ و $a \in \Domain{M_1}$ را بگیرید. اگر $a$ از پیش
  در دامنهٔ~$p$ باشد، $q=p$ را بگیرید. اگر $p=\emptyset$، هر
  $b \in \Domain{M_2}$ را برگزینید و
  $q=\{\langle a,b\rangle\}$ را بگیرید. در حالت باقی‌مانده، فرض کنید
  $p(a_i)=b_i$ برای $i = 1$، \dots،~$n$ و، بدون کاستن از کلیت،
  $a_1 <_1 a_2 <_1 \cdots <_1 a_n$. برای
  $a \in \Domain{M_1}$ که در دامنهٔ~$p$ نیست، $b \in \Domain{M_2}$ را
  چنین بیابید:
  \begin{enumerate}
  \item اگر $a <_1 a_1$، $b \in \Domain{M_2}$ را چنان بگیرید که
    $b <_2 b_1$؛
  \item اگر $a_n <_1 a$، $b \in \Domain{M_2}$ را چنان بگیرید که
    $b_n <_2 b$؛
  \item اگر برای یک $i$، $a_i <_1 a <_1 a_{i+1}$، آنگاه
    $b \in \Domain{M_2}$ را چنان بگیرید که
    $b_i <_2 b <_2 b_{i+1}$.
  \end{enumerate}
  چون $\Struct{M_2}$ یک ترتیب خطی چگال بدون نقاط انتهایی است، همواره
  می‌توان $b$ای با خاصیت مطلوب یافت. تعریف کنید
  $q = p \cup \{ \langle a, b \rangle \}$ تا $q \in \PIso{I}$ بسط
  مطلوب~$p$ باشد. این خاصیت پیش را برقرار می‌کند. خاصیت پس مشابه است.
  پس $\Struct{M_1} \iso[p] \Struct{M_2}$؛ بنا بر
  \olref[pis]{thm:p-isom1}، $\Struct{M_1} \iso \Struct{M_2}$.
\end{proof}

\begin{prob}
  برهان \olref[mod][bas][dlo]{thm:cantorQ} را با وارسی اینکه
  $\PIso{I}$ خاصیت پس را برآورده می‌کند کامل کنید.
\end{prob}

\begin{rem}
  بگذارید $\Struct{S}$ هر ترتیب خطی چگال بدون نقاط انتهایی و
  !!{enumerable} باشد. آنگاه (بنا بر \olref{thm:cantorQ})
  $\Struct{S} \iso \Struct{Q}$، که در آن
  $\Struct{Q} = (\Rat, <)$ ترتیب خطی چگال و !!{enumerable}ی است که
  مجموعهٔ~$\Rat$ از اعداد گویا دامنهٔ آن است. اکنون دوباره
  !!{structure}~$\Struct{R} = (\Real, <)$ از
  \olref[thm]{remark:R} را در نظر بگیرید. دیدیم که !!a{enumerable}
  !!{structure}~$\Struct{S}$ای وجود دارد که
  $\Struct{R} \elemequiv \Struct{S}$. اما $\Struct{S}$
  !!a{enumerable} ترتیب خطی چگال بدون نقاط انتهایی است، پس با
  !!{structure}~$\Struct{Q}$ همریخت (و در نتیجه هم‌ارز ابتدایی) است.
  بنا بر تعدی هم‌ارزی ابتدایی، $\Struct{R} \elemequiv \Struct{Q}$.
  (می‌توانستیم این را مستقیماً با برقراری
  $\Struct{R} \iso[p] \Struct{Q}$ از طریق همان استدلال پس‌وپیش نشان
  دهیم.)
\end{rem}
\end{document}
